\section{问题三讨论}

结果显示，成本与场景效用之间不存在可跨场景复用的单一排序。DeepSeek-
V4-Flash Max 在三个标准化工作负载中均提供最低成本，并具有稳定的
Pareto 前沿概率；但在 General 场景中，Qwen3.8-Max、Gemini-3.1-Pro
(High) 和 GPT-5.6 Sol 的效用增益足以使其进入前沿。相反，在 Coding 和
Research 中，较高成本模型的效用增益主要集中在从 DeepSeek-V4-Pro Max
转向 Kimi K3 的末端区间，该转移的 ICER 较大，说明预算决策应明确区分
“获得更高效用”与“为小幅增益支付多少成本”。

Bootstrap 结果支持上述结论的稳健部分，但不应被解释为所有现实部署条件
下的模型排名。该概率是在固定 FULL cohort 成本和 Q2 效用抽样下计算的，
没有覆盖价格更新、缓存命中率、峰谷时段、延迟、可用性、速率限制和部署
约束。成本--性能拟合也只是横截面描述；Coding 的全体模型线性和对数拟合
解释度很低，General 的饱和拟合虽有较高 $R^2$，但留一交叉验证误差很大，
且 Coding 与 Research 的 Pareto 前沿拟合仅有 3 个点，不能据此作稳定的
函数形式判断。

本轮最重要的发布限制是价格观测完整性。Claude Fable 5 的 base SKU 价格
不能替代带 fallback 的完整配置成本，GLM-5.2 也不能用相邻 SKU 或第三方
报价补齐。因此，当前结果适合进入论文的“方法与暂定结果”部分，最终推荐
仍需人工核对价格来源、配置含义和计费输出 token 定义后再发布。
